\documentclass[11pt]{article}
\usepackage{amsmath,amsthm,amssymb}
\usepackage[margin=1in]{geometry}
\usepackage{stmaryrd}
\usepackage[colorlinks=true,linkcolor=blue,citecolor=blue,urlcolor=blue]{hyperref}

\newtheorem{theorem}{Theorem}
\newtheorem{proposition}[theorem]{Proposition}
\newtheorem{corollary}[theorem]{Corollary}
\newtheorem{lemma}[theorem]{Lemma}
\theoremstyle{definition}
\newtheorem{definition}[theorem]{Definition}
\newtheorem{remark}[theorem]{Remark}

\newcommand{\C}{\mathcal{C}}
\newcommand{\Ob}{\mathrm{Ob}}
\newcommand{\id}{\mathrm{id}}
\newcommand{\Z}{\mathbb{Z}}
\newcommand{\Aut}{\mathrm{Aut}}
\newcommand{\Inn}{\mathrm{Inn}}
\newcommand{\Out}{\mathrm{Out}}
\newcommand{\gerbe}{\mathfrak{g}}
\newcommand{\abar}{\bar{\alpha}}
\newcommand{\B}{\mathrm{B}}
\newcommand{\Ad}{\mathrm{Ad}}

\title{The inventory floor needs the 2-cell datum:\\ $D1$ versus $D2$}
\author{MacBeth\thanks{Kodamai. A follow-up note to the band-boundary review; the companion
proof file \texttt{2026-10-12-inventory-floor-d1-d2-dichotomy} and the Lean witness
\texttt{2026-10-12-floor-d1-separation-s3} live in the author's \texttt{work-in-progress}
corpus, content commit \texttt{f19cfcb}. Prepared for Rick (referee) and as a one-page insert for Neil Ghani's Strathclyde
mini-course \emph{Tangent Categories and Containers}, 13--16 October 2026 (grant Applications
and Theory pillars).}\\ \small Kodamai}
\date{12 October 2026; revised 13 October 2026 (Rick 2026-10-13 referee MUSTs: $\alpha$-strict standing
hypothesis; Thm~1.3 ``centre-too-small'' and $D2\Rightarrow D1$ retracted; clause-(a)/Def~(D2) scoped)\\[2pt]
\small Prepared for Rick (referee) and Neil Ghani (PI); CC Robin Langer. Content commit \texttt{ccbf3b6}.}

\begin{document}
\maketitle

\begin{abstract}
Abstracting a supply chain to a small category $\C$ with fungible stock glued up to a
relabelling group $K=\Aut F$, the obstruction to a \emph{strict global inventory} is a
$K$-gerbe over the nerve $N(\C)$, staged following Giraud: a nonabelian band
$\beta\in H^1(N(\C);\Out K)$ must vanish before an abelian class
$[\gerbe]\in H^2(N(\C);Z(K))$ is defined. The ``floor corollary'' claimed that over the
band-trivial stratum a strict global inventory exists exactly when $[\gerbe]=0$. This is
false, and the reason is instructive: the claim silently fixes a notion of ``strict'' that
is one cohomological degree below $[\gerbe]$. Throughout we take the transport to be a
\emph{strict} functor (trivial coherence datum $k\equiv e$): this is the regime in which
the band-triviality criterion applies and the dichotomy is clean. We separate the two
honest notions. Retaining
the coherence (2-cell) datum, a $D2$ global inventory exists iff $[\gerbe]=0$ in
$H^2(N(\C);Z(K))$; discarding it, a $D1$ global inventory exists iff a degree-one class
$o_{D1}\in H^1(\C;\Inn K)/\mathrm{im}\,H^0(\C;\Out K)$ vanishes. The $D1$ condition implies
the $D2$ condition ($o_{D1}=0\Rightarrow[\gerbe]=0$), and the implication is strict; the
converse $D2\Rightarrow D1$ is \emph{false}. Already over a one-dimensional base, where
$H^2(N(\C);Z(K))=0$ forces $[\gerbe]=0$ for every $K$, the class $o_{D1}$ can still detect
a nontrivial inner transport. The witness is $\C=\B(\Z)$,
$K=S_3$, where $\beta=0$ and $[\gerbe]=0$ yet no naive strictification exists; its finite
core is machine-checked, and an independent sweep of $248$ strict-$\alpha$ cases
corroborates the dichotomy with the designed controls failing as they must.
\end{abstract}

\section{Introduction}\label{sec:intro}

\paragraph{The floor, and a false corollary.} A supply chain is a small category $\C$:
objects are nodes, morphisms are logistics operations, composition is sequencing. When
local inventory charts are comparable only \emph{up to a relabelling} of stock states --
up to the automorphism group $K:=\Aut F$ of the stock fibre $F$ -- the gluing datum is an
$\Aut F$-gerbe over the nerve $N(\C)$, and the obstruction to a \emph{strict global
inventory} (a single coherent trivialization across the chain) is staged in the manner of
Giraud~\cite{giraud}: a nonabelian band $\beta\in H^1(N(\C);\Out K)$ must vanish before an
abelian gerbe class $[\gerbe]\in H^2(N(\C);Z(K))$ is even defined. On the band-trivial
stratum $\beta=0$ it is tempting to close the account with a clean biconditional:
\[
\text{(floor corollary, \emph{false})}\qquad
\text{strict global inventory exists}\iff \beta=0\ \wedge\ [\gerbe]=0 .
\]
It is false. Rick, refereeing the band-boundary note, recomputed the obstruction for
$\C=\B(\Z)$ with $K=S_3$ and found $\beta=0$ and $[\gerbe]=0$ yet no strictification: the
reflection transport cannot be relabelled to the identity. This note explains \emph{why},
and in doing so turns the corollary into a pair of correct theorems.

\paragraph{Standing hypothesis (strict transport).} Throughout this note the transport
$\alpha\colon\C\to\B(\Aut K)$ is a \emph{strict} functor: its coherence datum is trivial,
$k_{g,f}\equiv e$. This is the regime in which the band-triviality criterion
(clause~(a), \S\ref{sec:prelim}) applies verbatim, so that $\beta=0$ yields a strict inner
transport $\abar$, and in which Definition~(D2) below may omit the comparison datum. For a
genuinely non-strict $\alpha$ (nontrivial $k_{g,f}$) neither holds as printed: the
band-trivial stratum need not yield a strict $\abar$ even when $\beta=0$ and $[\gerbe]=0$
(already for $\C=\B(\Z/3)$, or $K=D_4$), and both notions must then carry the $k$-datum.
That extension is deliberately outside the present scope; everything below is stated under
$k_{g,f}\equiv e$.

\paragraph{The gap: two meanings of ``strict''.} The corollary hides a choice. Once the
band is trivial, the transport is gauge-equivalent to a strict \emph{inner} functor
$\abar\colon\C\to\B(\Inn K)$ (Section~\ref{sec:prelim}); the question is what we ask of it.
Two answers:
\begin{itemize}
\item[\textbf{(D1)}] \emph{(naive)} a vertex gauge $\mu\colon\Ob\C\to\Aut K$ relabelling the
transport away \emph{entirely}: $\mu(y)\,\abar(f)\,\mu(x)^{-1}=\id_K$ for every $f\colon x\to y$.
No coherence datum is retained.
\item[\textbf{(D2)}] a gauge $\mu$ \emph{together with} a strict $K$-valued functor
$T\colon\C\to\B K$ ($T_{gf}=T_g T_f$) realising the transport by honest group elements:
$\mu(y)\,\abar(f)\,\mu(x)^{-1}=c_{T_f}$, where $c_a=\mathrm{conj}(a)$. (Under the standing
strict-transport hypothesis $\abar$ is strict, so this equation carries no coherence term;
for non-strict $\alpha$ it must be written with the comparison datum $k_{g,f}$, and then
$T$ need not be strict --- exactly the extension we exclude.)
\end{itemize}
$D1$ demands that the transport be \emph{trivialised}; $D2$ demands only that it \emph{lift}
to the group $K$, surviving as a genuine inner relabelling carried by elements of $K$. These
are different requests, and -- this is the whole point -- they live one cohomological degree
apart.

\paragraph{The result.} Write $Z=Z(K)$, $\Inn K=K/Z$, $\Out K=\Aut K/\Inn K$, $q\colon\Aut
K\twoheadrightarrow\Out K$, and let $[\gerbe]:=\abar^{*}[\star]\in H^2(N(\C);Z)$ be the
pullback along $\abar$ of the class $[\star]$ of the central extension $1\to Z\to K\to\Inn
K\to 1$ (Section~\ref{sec:prelim}; this is the gerbe class of~\cite{macbeth-capstone}).

\begin{theorem}[Floor dichotomy]\label{thm:main}
Fix a small category $\C$ with band-trivial transport $\abar\colon\C\to\B(\Inn K)$.
\begin{enumerate}
\item \textbf{($D2$ end.)} A $D2$ global inventory exists if and only if
$[\gerbe]=0\in H^2(N(\C);Z)$.
\item \textbf{($D1$ end.)} A $D1$ global inventory exists if and only if $o_{D1}=0$, where
$o_{D1}$ is the class of $\abar$ in the pointed quotient
$H^1(\C;\Inn K)/\mathrm{im}\,H^0(\C;\Out K)$. This is a \emph{degree-one} obstruction; it
is not $[\gerbe]$.
\item \textbf{($D1\Rightarrow D2$, strictly.)} $o_{D1}=0\Rightarrow[\gerbe]=0$, and the
implication is strict. The converse $D2\Rightarrow D1$ fails: already over a one-dimensional
base, where $H^2(N(\C);Z)=0$ forces $[\gerbe]=0$ for every $K$, the class $o_{D1}$ can be
nonzero (part~4). Centre size is not the governing parameter --- the gap occurs both over
centreless stock ($K=S_3$) and over stock with nontrivial centre ($K=D_4$, $Z=\Z/2$).
\item \textbf{(Separation.)} For $\C=\B(\Z)$, $K=S_3$, $\alpha=c_{(0\,1)}$: $\beta=0$ and
$[\gerbe]=0$ (so $D2$ holds -- an explicit lift is $T(n)=\sigma^n$), yet $o_{D1}\neq 0$
(no naive gauge strictifies the reflection). Hence the floor corollary is false, and the
honest $D1$ floor lives one degree down.
\end{enumerate}
\end{theorem}

\paragraph{The method.} One slogan runs through the proof: \emph{name the extension, and
its kernel names both the coefficient group and the cohomological degree.} $D2$ asks the
inner transport to lift through the central extension $(\star)$ with kernel $Z$; the
lifting obstruction is a degree-two class with coefficients in $Z$ -- exactly $[\gerbe]$
(Lemma~\ref{lem:lift}, the Eilenberg--Mac\,Lane / Brown criterion~\cite{eml,brown}). $D1$
asks the transport to be a \emph{coboundary} in the nonabelian $H^1(\C;\Inn K)$ -- a
degree-one condition with coefficients in $\Inn K$. Discarding the coherence datum costs
precisely one degree, and shifts the coefficient group from the centre $Z$ up to the inner
group $\Inn K$. Keeping these two coefficient groups distinct is the entire content of the
repair.

\paragraph{Consequences.} Over a one-dimensional base the floor is always clear on the
$D2$ side: $N(\B(\Z))\simeq S^1$ has $H^2=0$, so a $D2$ inventory exists for every $K$
(Corollary~\ref{cor:bz}). The $D1$ failure for $S_3$ is therefore not a failure of the
gerbe class at all -- it is the degree-one shadow the corollary could not see. For the
grant's Applications pillar the moral is operational: whether a supply chain admits a
coherent global relabelling depends on \emph{which} notion of ``global relabelling'' the
operator can afford. If the stock is genuinely fungible (centreless relabelling group),
demanding the strong form $D1$ can fail even when the weak form $D2$ -- carry the
relabelling as an explicit, composable group-valued correction -- succeeds.

\paragraph{Outline.} Section~\ref{sec:prelim} fixes the extension $(\star)$, the nonabelian
$H^1$, and the band-trivialised transport. Section~\ref{sec:d2} proves the $D2$ end
(Theorem~\ref{thm:main}.1). Section~\ref{sec:d1} constructs $o_{D1}$ and proves the $D1$
end (\ref{thm:main}.2). Section~\ref{sec:degree} records the degree lane and
$D1\Rightarrow D2$ (\ref{thm:main}.3). Section~\ref{sec:sep} works the $S_3$ separation
(\ref{thm:main}.4) and cites its machine-checked core.

\section{Preliminaries}\label{sec:prelim}

We recall only what the proof uses. Let $K=\Aut F$ be the relabelling group of the stock
fibre, $Z=Z(K)$ its centre, $\Inn K=K/Z$ the inner automorphism group, $\Out K=\Aut
K/\Inn K$ the outer group; $q\colon\Aut K\twoheadrightarrow\Out K$ is the projection and
$c\colon K\twoheadrightarrow\Inn K\subseteq\Aut K$, $c_a=\mathrm{conj}(a)$, is conjugation.

\paragraph{The central extension.} Conjugation gives a central extension of groups
\begin{equation}\tag{$\star$}\label{eq:star}
1\longrightarrow Z\longrightarrow K\xrightarrow{\ c\ }\Inn K\longrightarrow 1,
\qquad [\star]\in H^2(\Inn K;Z),
\end{equation}
central because inner automorphisms fix $Z$ pointwise, so the action is trivial.

\paragraph{The transport and its band.} The inventory datum is an \emph{outer transport}:
a pseudofunctor $\C\to\B(\Aut K)$ sending $f\colon x\to y$ to $\alpha(f)\in\Aut K$ with
coherence (2-cell) data $k_{g,f}\in K$, $\alpha(gf)=c_{k_{g,f}}\,\alpha(g)\alpha(f)$. Its
\emph{band} is the strict functor $\rho=q\circ\alpha\colon\C\to\B(\Out K)$, with class
$\beta=[\rho]\in H^1(N(\C);\Out K)$ (nonabelian $H^1$). Under the standing hypothesis
(\S\ref{sec:intro}) $\alpha$ is \emph{strict}, $k_{g,f}\equiv e$: an honest functor
$\C\to\B(\Aut K)$, with $\rho$ its band.

\paragraph{Standing hypothesis: the band is trivial.} Throughout, $\beta=0$. By the
band-triviality criterion (clause~(a) of the band-boundary note, computed) \emph{for strict
$\alpha$}, the functor $\alpha$ is $\Aut K$-gauge-equivalent to a \emph{strict} $\Inn
K$-valued functor. (Clause~(a) is stated and verified for strict $\alpha$; for non-strict
$\alpha$ it can fail even at $\beta=0$, which is precisely why strictness is a standing
hypothesis.) Fix such a gauge once and write $\abar\colon\C\to\B(\Inn K)$,
$\abar(gf)=\abar(g)\abar(f)$, $\abar(\id)=\id$, for the band-trivialised transport. This is
the object both $D1$ and $D2$ act on.

\paragraph{Nonabelian $H^1$.} For a group $G$ with trivial $\C$-action, let
$H^1(\C;G)$ be the pointed set of functors $\C\to\B G$ modulo natural isomorphism
(Serre~\cite[I.5]{serre}). Concretely a functor is a \emph{$1$-cocycle} $f\mapsto g_f\in G$
with $g_{\id}=e$ and $g_{gf}=g_g g_f$; a natural isomorphism is \emph{vertex data}
$\nu\colon\Ob\C\to G$ acting by $g_f\mapsto\nu(y)g_f\nu(x)^{-1}$. The basepoint $*$ is the
trivial cocycle $f\mapsto e$. A cocycle is a \emph{coboundary} (equivalent to $*$) iff
$g_f=\nu(y)^{-1}\nu(x)$ for some $\nu$. For $G$ nonabelian this is only a pointed set; that
is all we need. The transport $\abar$ is such a $1$-cocycle valued in $\Inn K$; for
$\C=\B(\Z)$, $H^1(\B(\Z);G)=G/\mathrm{conj}$ and $[\abar]$ is the conjugacy class of
$\abar(1)$.

\paragraph{The gerbe class.} Define $[\gerbe]:=\abar^{*}[\star]\in H^2(N(\C);Z)$, the
pullback of the extension class~\eqref{eq:star} along the strict functor
$\abar\colon N(\C)\to\B(\Inn K)$. That this is the $\Aut F$-gerbe's obstruction, and is
independent of the coherence datum, is the content of the capstone's rectification
theorem~\cite{macbeth-capstone}; we take it as given.

\section{The $D2$ floor: the obstruction is $[\gerbe]\in H^2$}\label{sec:d2}

The key idea is that $D2$ is exactly a \emph{strict lifting problem} through the central
extension~\eqref{eq:star}, and such problems are governed by a degree-two class with
coefficients in the kernel $Z$.

\begin{lemma}[Strict lifting through a central extension]\label{lem:lift}
Let $\varphi\colon\C\to\B(\Inn K)$ be a strict functor. A strict functor $T\colon\C\to\B K$
with $c\circ T=\varphi$ (that is, $T_f\in K$, $c(T_f)=\varphi_f$, $T_{gf}=T_g T_f$) exists if
and only if $\varphi^{*}[\star]=0\in H^2(N(\C);Z)$.
\end{lemma}

\begin{proof}
Choose a set-section $s\colon\Inn K\to K$ of $c$ with $s(e)=e$ and put $t_f:=s(\varphi_f)$,
so $c(t_f)=\varphi_f$. The defect $z(g,f):=t_{gf}(t_g t_f)^{-1}$ satisfies $c(z(g,f))=
\varphi_{gf}(\varphi_g\varphi_f)^{-1}=e$, hence $z(g,f)\in Z$. Associativity of composition
in $\C$ forces the $2$-cocycle identity $\delta z=0$ on the nerve $N(\C)$ (the standard bar
computation; $Z$ central, so the action is trivial), and $[z]\in H^2(N(\C);Z)$ is exactly
$\varphi^{*}[\star]$ -- the pullback of the bar $2$-cocycle of~\eqref{eq:star} along
$\varphi$~\cite[IV.3]{brown},~\cite{eml}. A strict lift is a choice $T_f=t_f\zeta_f$ with
$\zeta_f\in Z$ making the defect vanish; since $z'=z\cdot\delta\zeta$, a strict lift exists
iff $z$ is a $Z$-valued coboundary iff $[z]=0$.
\end{proof}

\begin{proposition}[$D2$ end]\label{prop:d2}
A $D2$ global inventory exists if and only if $[\gerbe]=\abar^{*}[\star]=0\in H^2(N(\C);Z)$.
\end{proposition}

\begin{proof}
Given $(\mu,T)$, set $\tau(f):=\mu(y)\,\abar(f)\,\mu(x)^{-1}$. Since $\Inn K\trianglelefteq
\Aut K$ and $\abar(f)\in\Inn K$, $\tau(f)$ is inner, and strictness of $\abar$ gives
$\tau(gf)=\mu(z)\abar(gf)\mu(x)^{-1}=\tau(g)\tau(f)$, so $\tau\colon\C\to\B(\Inn K)$ is a
\emph{strict} functor. The $D2$ equation $\tau(f)=c_{T_f}$ with $T$ strict says precisely
that $T$ is a strict lift of $\tau$ through~\eqref{eq:star}; by Lemma~\ref{lem:lift} such
$(\mu,T)$ exists iff $\tau^{*}[\star]=0$ for some admissible $\mu$.

It remains to see $\tau^{*}[\star]=0\iff\abar^{*}[\star]=0$. Decompose $\mu(x)=\hat h_c\,j(x)$
on each component $c$, where $\hat h_c\in\Aut K$ is a fixed lift of the locally constant
outer part $q\circ\mu$ and $j(x)\in\Inn K$. Then
$\tau(f)=\Ad(\hat h_c)\big(j(y)\,\abar(f)\,j(x)^{-1}\big)=\Ad(\hat h_c)\circ\abar'(f)$, where
$\abar'(f):=j(y)\abar(f)j(x)^{-1}$ is naturally isomorphic to $\abar$ (inner vertex data), so
$\abar'^{*}=\abar^{*}$ on $H^2$. The global automorphism $\hat h_c$ preserves~\eqref{eq:star}
(it fixes $Z$ setwise, inducing $\hat h_c|_Z\in\Aut Z$, and covers $\Ad(\hat h_c)\in\Aut(\Inn
K)$), so $\tau^{*}[\star]=(\hat h_c|_Z)_{*}\,\abar^{*}[\star]=(\hat h_c|_Z)_{*}[\gerbe]$, an
automorphism image of $[\gerbe]$. An automorphism of $H^2(N(\C);Z)$ fixes $0$, so
$\tau^{*}[\star]=0\iff[\gerbe]=0$.
\end{proof}

\begin{corollary}\label{cor:bz}
For $\C=\B(\Z)$, $N(\C)\simeq S^1$, so $H^2(S^1;Z)=0$ and $[\gerbe]=0$ for every $K$: a $D2$
global inventory always exists over a free one-generator base.
\end{corollary}

\section{The $D1$ floor: the obstruction drops to $o_{D1}\in H^1$}\label{sec:d1}

The key idea here is dual to the last: $D1$ asks $\abar$ not to \emph{lift} but to
\emph{vanish} -- to be a coboundary in the nonabelian $H^1(\C;\Inn K)$. This is a
degree-one condition, and its coefficient group is $\Inn K$, not $Z$.

\paragraph{Band-preservation is forced.} If a $D1$ gauge $\mu$ exists, applying $q$ to
$\mu(y)\abar(f)\mu(x)^{-1}=\id$ and using $q(\abar(f))=e$ (band trivial) gives
$q(\mu(y))=q(\mu(x))$ for every $f\colon x\to y$. Hence $q\circ\mu\colon\Ob\C\to\Out K$ is
locally constant: $D1$ \emph{entails} band-preservation, with no extra hypothesis.

\paragraph{Construction of $o_{D1}$.} $H^0(\C;\Out K)$ is the group of locally constant
maps $\Ob\C\to\Out K$ ($=\Out K$ per component). It acts on the pointed set
$H^1(\C;\Inn K)$: given $\theta\in H^0(\C;\Out K)$, choose a locally constant lift $\hat
h\colon\Ob\C\to\Aut K$ with $q\circ\hat h=\theta$, and send a cocycle $g_f$ to $\hat
h(y)\,g_f\,\hat h(x)^{-1}\in\Inn K$ (inner, as $\Inn K\trianglelefteq\Aut K$). A different
lift of $\theta$ differs by an $\Inn K$-valued cochain, changing the result by an $\Inn
K$-coboundary; so the action descends to $H^1(\C;\Inn K)$ and fixes the basepoint $*$. Set
\[
o_{D1}:=[\abar]\ \in\ H^1(\C;\Inn K)\,/\,\mathrm{im}\,H^0(\C;\Out K).
\]

\begin{proposition}[$D1$ end]\label{prop:d1}
A $D1$ global inventory exists if and only if $o_{D1}=0$; equivalently, iff $[\abar]=*$ in
$H^1(\C;\Inn K)$.
\end{proposition}

\begin{proof}
We first show $D1\iff[\abar]=*$.
\begin{itemize}
\item If $[\abar]=*$, pick $j\colon\Ob\C\to\Inn K$ with $\abar(f)=j(y)^{-1}j(x)$. Set
$\mu:=j$. Then $\mu(y)\abar(f)\mu(x)^{-1}=j(y)\,j(y)^{-1}j(x)\,j(x)^{-1}=\id$, so $D1$ holds.
\item Conversely let $\mu\colon\Ob\C\to\Aut K$ solve $D1$. By the forced constraint
$q\circ\mu=\theta_c$ is constant on each component $c$; choose a lift $\hat h_c\in\Aut K$ of
$\theta_c$ and write $\mu(x)=\hat h_c\,j(x)$ with $j(x):=\hat h_c^{-1}\mu(x)\in\ker q=\Inn
K$. For $f\colon x\to y$ in $c$, $\abar(f)=\mu(y)^{-1}\mu(x)=j(y)^{-1}j(x)$, an $\Inn
K$-coboundary on $c$. Over all components, $[\abar]=*$.
\end{itemize}
Now the quotient. The $\Out K$-action fixes $*$, so the fibre of the quotient map over the
basepoint is exactly the $\Out K$-orbit of $*$, namely $\{*\}$. Hence $o_{D1}=0$ (the
basepoint of the quotient) iff $[\abar]=*$ iff $D1$ holds.
\end{proof}

The quotient by $\mathrm{im}\,H^0(\C;\Out K)$ makes $o_{D1}$ independent of the choice of
fibre coordinates (a global automorphism of $K$ acts by an outer class). The action can be
genuinely nontrivial -- for $K=D_4$ it swaps the two reflection conjugacy classes, both
nonzero -- so the quotient is needed for coordinate-freeness, but by the last line of the
proof it never changes \emph{whether} $o_{D1}=0$: the quotient is vanishing-neutral.

\section{The degree lane, and $D1\Rightarrow D2$}\label{sec:degree}

\begin{remark}[one well-definedness caveat, and the degree drop]\label{rem:degree}
Both invariants are canonical only up to the residual $\Out K$-action (global fibre
re-coordinatisation): $[\gerbe]$ up to the $\Out K$-action on $H^2(N(\C);Z)$
(Proposition~\ref{prop:d2}), $o_{D1}$ by the explicit quotient (Proposition~\ref{prop:d1}).
In both cases the action is by automorphisms fixing $0$ (resp.\ $*$), so vanishing is
invariant and both biconditionals are well-posed. $D1$ and $D2$ differ not in their gauge
symmetry but in their \emph{degree}: $D1$ asks $\abar$ to be a coboundary -- trivialised to
the identity, a degree-one ($H^1$) condition on $\Inn K$; $D2$ asks $\abar$ only to lift to
a strict $K$-valued functor, a degree-two ($H^2$, $k$-invariant) condition on $Z$. Dropping
the coherence datum costs exactly one cohomological degree and raises the coefficient group
from $Z$ to $\Inn K$.
\end{remark}

\begin{proposition}[$D1\Rightarrow D2$, strictly]\label{prop:implies}
$o_{D1}=0\Rightarrow[\gerbe]=0$, and the containment is strict.
\end{proposition}

\begin{proof}
If $[\abar]=*$ then $\abar$ is naturally isomorphic to the constant functor at the identity
object of $\B(\Inn K)$; pullback along a functor isomorphic to a constant kills $H^2$, so
$[\gerbe]=\abar^{*}[\star]=0$ (equivalently, take $T\equiv e$ in $D2$). Strictness: the
converse $D2\Rightarrow D1$ fails already over a one-dimensional base, where
$H^2(N(\C);Z)=0$ forces $[\gerbe]=0$ for every $K$ while $\abar$ may be a nontrivial inner
transport. $\C=\B(\Z)$, $K=S_3$ (\S\ref{sec:sep}) is such a witness, with $Z=1$; so is
$K=D_4$ with $\abar=c_s$ a reflection, where $Z=\Z/2\neq 1$ --- so centre size does not
govern the gap, and the earlier ``$Z$ too small'' heuristic is not a characterisation.
\end{proof}

\begin{remark}[what the gap is, and is not]\label{rem:notcentre}
The gap $o_{D1}\neq 0=[\gerbe]$ is a phenomenon of cohomological \emph{degree}, not of
centre size. On a base with $H^2=0$ the degree-two class $[\gerbe]$ is forced to vanish for
\emph{every} coefficient group, so $D2$ holds unconditionally, yet the degree-one class
$o_{D1}$ still records the inner transport. An exact description of the locus
$\{\,o_{D1}\neq 0=[\gerbe]\,\}$ over a general base is not attempted here.
\end{remark}

\section{The separation: $\C=\B(\Z)$, $K=S_3$}\label{sec:sep}

$S_3$ is \emph{complete}: $Z(S_3)=1$, $\Inn S_3=S_3$, $\Out S_3=1$, $\Aut S_3\cong
S_3$~\cite[Thm 7.5]{rotman}. The transport on $\B(\Z)$ is determined by the generator,
$\alpha(1)=c_\sigma$ with $\sigma=(0\,1)$, and freely $\alpha(n)=c_{\sigma^n}$; as $\B(\Z)$
is free on one generator this is already strict with trivial coherence datum, so
$\abar(n)=c_{\sigma^n}$.

\begin{itemize}
\item \textbf{Band.} $\Out S_3=1$, so $\beta\in H^1(S^1;1)=0$. \checkmark
\item \textbf{Gerbe class.} $Z(S_3)=1$, so $H^2(S^1;Z)=0$ and $[\gerbe]=0$; independently,
$N(\B(\Z))\simeq S^1$ forces $H^2=0$ (Corollary~\ref{cor:bz}). By
Proposition~\ref{prop:d2} a $D2$ inventory exists, with explicit lift $T(n)=\sigma^n$: a
homomorphism $\Z\to S_3$ since $\sigma^2=e$, and $\mu=\id$, $c_{T(n)}=c_{\sigma^n}=\abar(n)$.
\checkmark
\item \textbf{$o_{D1}$.} $H^1(\B(\Z);\Inn S_3)=S_3/\mathrm{conj}=\{[e],[\text{transposition}],
[\text{3-cycle}]\}$, and $o_{D1}$ is the class of $\sigma$, a transposition $\neq[e]$; with
$\Out=1$ it is left unchanged, so $o_{D1}\neq 0$. By Proposition~\ref{prop:d1} no $D1$
inventory exists. Directly: on one object $\mu\,\abar(1)\,\mu^{-1}=c_{m\sigma m^{-1}}$, and
$c_{m\sigma m^{-1}}=\id\iff m\sigma m^{-1}=e$ (as $Z(S_3)=1$ makes $c$ injective),
impossible for $\sigma\neq e$. \checkmark
\end{itemize}

Hence $\beta=0\wedge[\gerbe]=0$ while $D1$ fails: the floor corollary is false, and the
honest $D1$ floor lives one degree down in $H^1$. The finite, decidable core of this
witness -- $Z(S_3)=1$, $\Aut S_3$ of order $6$ (so $\Out S_3=1$), the $D2$ lift $\sigma^2=e$,
and the exhaustive check that no $m\in S_3$ conjugates $\sigma$ to $e$ -- is machine-checked
in Lean~4 (\texttt{floor\_separation}, kernel-checked with only \texttt{propext} and
\texttt{Quot.sound}, no \texttt{native\_decide} and no added axioms). The cohomological
identifications ($[\gerbe]=\abar^{*}[\star]$, the location of $o_{D1}$, and the forcing of
vanishing from trivial coefficient groups) are the cited classical bridges, not re-derived
in Lean. Beyond this witness, an independent recomputation of the dichotomy of
Theorem~\ref{thm:main} over $248$ strict-$\alpha$ cases (independent referee recomputation;
\texttt{work-in-progress scripts/day238}) found $0$ violations, with the designed controls
failing exactly as they must.

\section*{Upshot}

Rick's $B(\Z)/S_3$ counterexample is now a theorem. The inventory floor splits cleanly by
\emph{how much structure one retains}. Keep the coherence (2-cell) datum ($D2$): the
obstruction is the gerbe class $[\gerbe]\in H^2(N(\C);Z(K))$. Discard it ($D1$): the
obstruction drops a degree to $o_{D1}\in H^1(\C;\Inn K)/\mathrm{im}\,H^0(\C;\Out K)$. The
$D1$ condition implies the $D2$ condition, strictly; the converse fails already over a
one-dimensional base, where $H^2=0$ makes $[\gerbe]$ vanish for every $K$ while $o_{D1}$
still sees a nontrivial inner transport. Both $\B(\Z)/S_3$ (centreless) and $\B(\Z)/D_4$
(centre $\Z/2$) are honest witnesses: the gap is a matter of cohomological degree, not of
centre size. The repair isolates the single recurring
hazard in the program -- confusing the central coefficient $Z(K)$ of the degree-two class
with the inner coefficient $\Inn K$ of the degree-one one. Whether this re-promotes the
band-boundary note is a separate registry judgement, deferred; here we have stated only
what is proved. The scope is the band-trivial stratum $\beta=0$; the nontrivial-band and
$T^2$-coupled frontiers are treated in the companion capstone~\cite{macbeth-capstone}.

\begin{thebibliography}{9}
\bibitem{brown} K.~S.~Brown, \emph{Cohomology of Groups}, Graduate Texts in Mathematics 87,
Springer, 1982 (Chapter IV.3).
\bibitem{eml} S.~Eilenberg and S.~Mac\,Lane, \emph{Cohomology theory in abstract groups.~II.
Group extensions with a non-Abelian kernel}, Ann.\ of Math.\ (2) \textbf{48} (1947), 326--341.
\bibitem{giraud} J.~Giraud, \emph{Cohomologie non abélienne}, Grundlehren der mathematischen
Wissenschaften 179, Springer, 1971 (Chapters III--IV).
\bibitem{macbeth-capstone} MacBeth, \emph{The inventory gerbe couples on the torus}, Kodamai
work-in-progress note, content commit \texttt{e306ae5}, 2026.
\bibitem{rotman} J.~J.~Rotman, \emph{An Introduction to the Theory of Groups}, 4th ed.,
Graduate Texts in Mathematics 148, Springer, 1995 (Theorem 7.5).
\bibitem{serre} J.-P.~Serre, \emph{Galois Cohomology}, Springer Monographs in Mathematics,
Springer, 1997 (Chapter I.5).
\end{thebibliography}

\end{document}
